% Unit 061 English translation of chapter4.tex lines 2116--2267.
% Source: Methods of Algebra, Volume 2, commit
% 9a5803ff77dd3257484cb177f851a73770a59dd3 (CC BY 4.0).
% stable-unit-id: o014.aljabr2.chapter4.unbounded-derived-functors
% Coverage: complete Section 4.11; stops before Section 4.12 at line 2268.

% segment-id: o014.li2.chapter4.unbounded-derived-functors.h001
\section{Unbounded Derived Functors}\label{sec:unbounded-derived}
\hypertarget{o014.aljabr2.chapter4.unbounded-derived-functors}{ }
% segment-id: o014.li2.chapter4.unbounded-derived-functors.x001
\index{derived functor!unbounded}

% segment-id: o014.li2.chapter4.unbounded-derived-functors.p001
Consider an additive functor $F:\mathcal{A}\to\mathcal{A}'$ between abelian
categories. This section returns to Definition~\ref{def:derived-functor} and
focuses on the versions of derived functors for unbounded derived categories,
% segment-id: o014.li2.chapter4.unbounded-derived-functors.q001
\[ \mathrm{R}F, \; \mathrm{L}F: \cate{D}(\mathcal{A}) \to \cate{D}(\mathcal{A}'). \]

% segment-id: o014.li2.chapter4.unbounded-derived-functors.p002
To guarantee the existence of $\mathrm{R}F$ or $\mathrm{L}F$ and to compute
with it, we need the notions of K-injective and K-projective complexes from
Definition~\ref{def:K-resolutions}, together with the corresponding
resolutions.

% segment-id: o014.li2.chapter4.unbounded-derived-functors.le001
\begin{lemma}\label{prop:K-injective-zero}
	Let $\mathcal{A}$ be an abelian category. If
	$Z\in\Obj(\cate{K}(\mathcal{A}))$ is both K-injective (respectively,
	K-projective) and acyclic, then $Z=0$.
\end{lemma}
% segment-id: o014.li2.chapter4.unbounded-derived-functors.pf001
\begin{proof}
	Lemma~\ref{prop:K-injective-criterion} implies that
	$\Hom_{\cate{K}(\mathcal{A})}(Z,Z)=0$.
\end{proof}

% segment-id: o014.li2.chapter4.unbounded-derived-functors.p003
We first state a general result. Fix triangulated categories
$\mathcal{D}_1$, $\mathcal{D}_2$, and $\mathcal{D}'$, together with their
saturated triangulated subcategories $\mathcal{N}_1$, $\mathcal{N}_2$, and
$\mathcal{N}'$.

% segment-id: o014.li2.chapter4.unbounded-derived-functors.le002
\begin{lemma}\label{prop:K-cats}
	Given a triangulated bifunctor
	$G:\mathcal{D}_1\times\mathcal{D}_2\to\mathcal{D}'$, define a full
	subcategory $\mathcal{K}_2^G$ of $\mathcal{D}_2$ by
% segment-id: o014.li2.chapter4.unbounded-derived-functors.q002
	\[ \Obj(\mathcal{K}_2^G) = \left\{ X_2 \in \Obj(\mathcal{D}_2) : X_1 \in \Obj(\mathcal{N}_1) \implies G(X_1, X_2) \in \Obj(\mathcal{N}') \right\}. \]
	Then $\mathcal{K}_2^G$ is a saturated triangulated subcategory. After
	interchanging the two variables, the corresponding statement holds for
	$\mathcal{K}_1^G\subset\mathcal{D}_1$.
\end{lemma}
% segment-id: o014.li2.chapter4.unbounded-derived-functors.pf002
\begin{proof}
	The isomorphism $G(X_1,X_2[1])\simeq G(X_1,X_2)[1]$ shows that
	$\mathcal{K}_2^G$ is closed under translation. Similarly, suppose that
	$I\to J\to K\xrightarrow{+1}$ is a distinguished triangle with
	$I,K\in\Obj(\mathcal{K}_2^G)$. From the distinguished triangle
% segment-id: o014.li2.chapter4.unbounded-derived-functors.q003
	\[ G(X_1, I) \to G(X_1, J) \to G(X_1, K) \xrightarrow{+1} \]
	it follows immediately that
	$X_1\in\Obj(\mathcal{N}_1)\implies
	G(X_1,J)\in\Obj(\mathcal{N}')$. Saturation is clear.
\end{proof}

% segment-id: o014.li2.chapter4.unbounded-derived-functors.p004
This result can be applied to
$\Hom:\mathcal{A}^{\opp}\times\mathcal{A}\to\cate{Ab}$ and the
corresponding triangulated bifunctor $\cate{K}_{\Pi}\Hom$.
Example~\ref{eg:Hom-bicplx} has shown that the latter bifunctor is essentially
the $\Hom$ complex:
% segment-id: o014.li2.chapter4.unbounded-derived-functors.e001
\begin{equation}\label{eqn:CPiHom}\begin{gathered}
	(\cate{C}_{\Pi} \Hom)(\sigma^{-1} X_1, X_2) \simeq \Hom^\bullet\left( X_1, X_2 \right), \\
	(\cate{K}_\Pi \Hom)(\sigma^{-1} X_1, X_2) \simeq \Hom^\bullet\left( X_1, X_2 \right)\; \text{viewed in $\cate{K}(\cate{Ab})$},
\end{gathered}\end{equation}
% segment-id: o014.li2.chapter4.unbounded-derived-functors.p005
where $\sigma:\cate{C}(\mathcal{A}^{\opp})\rightiso
\cate{C}(\mathcal{A})^{\opp}$ and
$\cate{K}(\mathcal{A}^{\opp})\rightiso\cate{K}(\mathcal{A})^{\opp}$ are
the isomorphisms from Definition--Proposition~\ref{def:sigma} and
Proposition~\ref{prop:sigma-homotopy}. We shall later need the corresponding
version for derived categories.

% segment-id: o014.li2.chapter4.unbounded-derived-functors.pr001
\begin{proposition}\label{prop:K-injective-F-injective}
	All K-injective (respectively, K-projective) complexes form a saturated
	triangulated subcategory $\mathcal{I}$ (respectively, $\mathcal{P}$) of
	$\cate{K}(\mathcal{A})$. If $\mathcal{A}$ has enough K-injective
	(respectively, K-projective) complexes, then for every additive functor
	$F:\mathcal{A}\to\mathcal{A}'$, the subcategory $\mathcal{I}$
	(respectively, $\mathcal{P}$) is $F$-injective (respectively,
	$F$-projective).
\end{proposition}
% segment-id: o014.li2.chapter4.unbounded-derived-functors.pf003
\begin{proof}
	It suffices to treat the K-injective case. The first assertion follows by
	applying Lemma~\ref{prop:K-cats} to $G=\cate{K}_{\Pi}\Hom$, using the
	definition of K-injective complexes in
	Definition~\ref{def:K-resolutions} and \eqref{eqn:CPiHom}: the subcategory
	$\mathcal{K}_2^G$ there is precisely $\mathcal{I}$ here.

	Suppose that $\mathcal{A}$ has enough K-injective complexes. If $I$ is an
	acyclic K-injective complex, Lemma~\ref{prop:K-injective-zero} implies that
	$I=0$, so $(\cate{K}F)I$ is certainly acyclic. Thus all the conditions in
	Definition~\ref{def:F-injective} for an $F$-injective triangulated
	subcategory are satisfied.
\end{proof}

% segment-id: o014.li2.chapter4.unbounded-derived-functors.p006
Write $Q:\cate{K}(\mathcal{A})\to\cate{D}(\mathcal{A})$ and
$Q':\cate{K}(\mathcal{A}')\to\cate{D}(\mathcal{A}')$ for the localization
functors.

% segment-id: o014.li2.chapter4.unbounded-derived-functors.co001
\begin{corollary}
	If $\mathcal{A}$ has enough K-injective complexes, then $F$ has a right
	derived functor
	$\mathrm{R}F:\cate{D}(\mathcal{A})\to\cate{D}(\mathcal{A}')$, such that
	$\mathrm{R}F(QX)\simeq Q'\cate{K}F(X)$ whenever $X$ is K-injective.

	Dually, if $\mathcal{A}$ has enough K-projective complexes, then $F$ has a
	left derived functor
	$\mathrm{L}F:\cate{D}(\mathcal{A})\to\cate{D}(\mathcal{A}')$, such that
	$\mathrm{L}F(QX)\simeq Q'\cate{K}F(X)$ whenever $X$ is K-projective.
\end{corollary}
% segment-id: o014.li2.chapter4.unbounded-derived-functors.pf004
\begin{proof}
	Apply Proposition~\ref{prop:K-injective-F-injective} in the general
	framework of Proposition~\ref{prop:localization-injective}.
\end{proof}

% segment-id: o014.li2.chapter4.unbounded-derived-functors.p007
Now fix abelian categories $\mathcal{A}_1$, $\mathcal{A}_2$, and
$\mathcal{B}$. For an additive bifunctor
$F:\mathcal{A}_1\times\mathcal{A}_2\to\mathcal{B}$, consider the unbounded
version of the derived bifunctor $\mathrm{R}F$ (respectively,
$\mathrm{L}F$) from Definition~\ref{def:derived-bifunctor}. From now on, we
always assume that $\mathcal{B}$ has countable products (respectively,
countable coproducts).

% segment-id: o014.li2.chapter4.unbounded-derived-functors.p008
Concretely, we wish to determine $\mathrm{R}F$ (respectively,
$\mathrm{L}F$) using the subcategories $\mathcal{I}_1$, $\mathcal{I}_2$
(respectively, $\mathcal{P}_1$, $\mathcal{P}_2$) above.

% segment-id: o014.li2.chapter4.unbounded-derived-functors.le003
\begin{lemma}\label{prop:bifunctor-via-K-injective-prep}
	For $i=1,2$, take the triangulated subcategory $\mathcal{I}_i$
	(respectively, $\mathcal{P}_i$) of $\cate{K}(\mathcal{A}_i)$ from
	Proposition~\ref{prop:K-injective-F-injective}.
% segment-id: o014.li2.chapter4.unbounded-derived-functors.l001
	\begin{compactitem}
% segment-id: o014.li2.chapter4.unbounded-derived-functors.i001
		\item If $\mathcal{A}_1$ and $\mathcal{A}_2$ have enough K-injective
		complexes, then $\mathcal{I}_1\times\mathcal{I}_2$ is $F$-injective.
% segment-id: o014.li2.chapter4.unbounded-derived-functors.i002
		\item If $\mathcal{A}_1$ and $\mathcal{A}_2$ have enough K-projective
		complexes, then $\mathcal{P}_1\times\mathcal{P}_2$ is $F$-projective.
	\end{compactitem}
\end{lemma}
% segment-id: o014.li2.chapter4.unbounded-derived-functors.pf005
\begin{proof}
	It suffices to treat the K-injective case. The key is to show that if
	$X_1\in\Obj(\mathcal{I}_1)$, then
	$(\cate{K}_{\Pi}F)(X_1,\cdot)$ sends every acyclic object $Z$ of
	$\mathcal{I}_2$ to an acyclic object of $\cate{K}(\mathcal{B})$. But
	Lemma~\ref{prop:K-injective-zero} gives $Z=0$, and
	$(\cate{K}_{\Pi}F)(X_1,\cdot)$ is additive, so the assertion follows
	immediately. The corresponding result follows after interchanging the
	indices $1$ and $2$.
\end{proof}

% segment-id: o014.li2.chapter4.unbounded-derived-functors.p009
Write $Q:\cate{K}(\mathcal{B})\to\cate{D}(\mathcal{B})$ and
$Q_i:\cate{K}(\mathcal{A}_i)\to\cate{D}(\mathcal{A}_i)$ for the
localization functors ($i=1,2$).

% segment-id: o014.li2.chapter4.unbounded-derived-functors.pr002
\begin{proposition}\label{prop:bifunctor-via-K-injective}
	For $F:\mathcal{A}_1\times\mathcal{A}_2\to\mathcal{B}$, suppose that
	$\mathcal{A}_1$ and $\mathcal{A}_2$ both have enough K-injective
	(respectively, K-projective) complexes. Then the derived bifunctor
	$\mathrm{R}F$ (respectively, $\mathrm{L}F$) exists. If
	$X_i\in\Obj(\cate{K}(\mathcal{A}))$ is K-injective (respectively,
	K-projective) for $i=1,2$, then there is a canonical isomorphism
% segment-id: o014.li2.chapter4.unbounded-derived-functors.q004
	\[ \mathrm{R}F(Q_1 X_1, Q_2 X_2) \simeq Q (\cate{K}_{\Pi} F)(X_1, X_2) \quad \text{(respectively, $\mathrm{L}F(Q_1 X_1, Q_2 X_2) \simeq Q (\cate{K}_{\oplus} F)(X_1, X_2)$) }. \]
\end{proposition}
% segment-id: o014.li2.chapter4.unbounded-derived-functors.pf006
\begin{proof}
	Apply Lemma~\ref{prop:bifunctor-via-K-injective-prep} directly to
	Proposition~\ref{prop:localization-injective-bifunctor}.
\end{proof}

% segment-id: o014.li2.chapter4.unbounded-derived-functors.p010
For a particular bifunctor $F$, one can often choose $F$-injective
(respectively, $F$-projective) subcategories larger than the subcategories of
K-injective (respectively, K-projective) complexes. The $\RHom$ functor to be
discussed next is one example.

% segment-id: o014.li2.chapter4.unbounded-derived-functors.th001
\begin{theorem}[Unbounded $\RHom$]\label{prop:unbounded-RHom}
% segment-id: o014.li2.chapter4.unbounded-derived-functors.x002
	\index{$\RHom$ functor}
% segment-id: o014.li2.chapter4.unbounded-derived-functors.x003
	\index[sym1]{RHom@$\RHom$}
	Suppose that the abelian category $\mathcal{A}$ has enough K-injective or
	enough K-projective complexes. For simplicity, denote every localization
	functor by $Q$.
% segment-id: o014.li2.chapter4.unbounded-derived-functors.l002
	\begin{enumerate}[(i)]
% segment-id: o014.li2.chapter4.unbounded-derived-functors.i003
		\item The bifunctor
		$\Hom=\Hom_{\mathcal{A}}:\mathcal{A}^{\opp}\times\mathcal{A}
		\to\cate{Ab}$ has a right derived bifunctor
		$\cate{D}(\mathcal{A}^{\opp})\times\cate{D}(\mathcal{A})
		\to\cate{D}(\cate{Ab})$. Via the equivalence $\sigma$ of
		Proposition~\ref{prop:derived-cat-op}, it can also be written as
% segment-id: o014.li2.chapter4.unbounded-derived-functors.q005
		\[ \RHom: \cate{D}(\mathcal{A})^{\opp} \times \cate{D}(\mathcal{A}) \to \cate{D}(\cate{Ab}). \]
		Moreover, if $\mathcal{A}$ has enough K-injective (respectively,
		K-projective) complexes, then
		$\cate{K}(\mathcal{A})^{\opp}\times\mathcal{I}$ (respectively,
		$\mathcal{P}^{\opp}\times\cate{K}(\mathcal{A})$) is
		$\Hom$-injective.
% segment-id: o014.li2.chapter4.unbounded-derived-functors.i004
		\item Suppose that $\mathcal{A}$ has enough K-injective complexes and
		that $X_2\to I_2$ is a K-injective resolution. This resolution induces
% segment-id: o014.li2.chapter4.unbounded-derived-functors.q006
		\[ \RHom(Q X_1, Q X_2) \rightiso Q \Hom^\bullet(X_1, I_2). \]
		In particular,
		$\RHom(QX_1,\cdot)\simeq
		\mathrm{R}\left(\Hom(X_1,\cdot)\right)$.
% segment-id: o014.li2.chapter4.unbounded-derived-functors.i005
		\item Suppose that $\mathcal{A}$ has enough K-projective complexes and
		that $P_1\to X_1$ is a K-projective resolution. This resolution induces
% segment-id: o014.li2.chapter4.unbounded-derived-functors.q007
		\[ \RHom(Q X_1, Q X_2) \rightiso Q \Hom^\bullet(P_1, X_2). \]
		In particular,
		$\RHom(\cdot,QX_2)\simeq
		\mathrm{R}\left(\Hom(\cdot,X_2)\right)$.
% segment-id: o014.li2.chapter4.unbounded-derived-functors.i006
		\item There is a canonical isomorphism
		$\Hm^n\RHom(X,Y)\simeq
		\Hom_{\cate{D}(\mathcal{A})}(X,Y[n])=: \Ext^n(X,Y)$
		(Definition~\ref{def:Ext-general}).
	\end{enumerate}
\end{theorem}
% segment-id: o014.li2.chapter4.unbounded-derived-functors.pf007
\begin{proof}
	The argument is nearly identical to that for
	Theorem~\ref{prop:RHom-bounded}. For example, suppose that $\mathcal{A}$
	has enough K-injective complexes. The key is to show that
	$\cate{K}(\mathcal{A})^{\opp}\times\mathcal{I}$ is $\Hom$-injective.
% segment-id: o014.li2.chapter4.unbounded-derived-functors.l003
	\begin{itemize}
% segment-id: o014.li2.chapter4.unbounded-derived-functors.i007
		\item Fix a complex $X_1$. The functor
		$(\cate{K}_\Pi\Hom)(X_1,\cdot)$ sends every acyclic K-injective complex
		$X_2$ to an acyclic complex, because $X_2$ is $0$ in
		$\cate{K}(\mathcal{A})$ (Lemma~\ref{prop:K-injective-zero}).
% segment-id: o014.li2.chapter4.unbounded-derived-functors.i008
		\item Fix a K-injective complex $X_2$. If $X_1$ is acyclic, the
		definition of a K-injective complex implies that
		$\Hom^\bullet\left(X_1,X_2\right)$ is acyclic. By
		\eqref{eqn:CPiHom}, $(\cate{K}_\Pi\Hom)(X_1,X_2)$ is therefore
		acyclic as well.
	\end{itemize}
	Thus the acyclicity conditions required for the $\Hom$-injective property
	have been verified, while the resolution conditions hold directly.
\end{proof}

% segment-id: o014.li2.chapter4.unbounded-derived-functors.th002
\begin{theorem}\label{prop:RHom-adjoint}
	Consider an adjoint pair of additive functors between abelian categories
% segment-id: o014.li2.chapter4.unbounded-derived-functors.g001
	$\begin{tikzcd}
		F: \mathcal{A} \arrow[r, shift left] & \mathcal{A}' \arrow[l, shift left] : G
	\end{tikzcd}$.
	Suppose that $\mathcal{A}$ has enough K-projective complexes and
	$\mathcal{A}'$ has enough K-injective complexes. Then there is a canonical
	isomorphism in $\cate{D}(\cate{Ab})$
% segment-id: o014.li2.chapter4.unbounded-derived-functors.q008
	\[ \RHom_{\mathcal{A'}}\left(\mathrm{L}F(X), Y\right) \simeq \RHom_{\mathcal{A}}\left(X, \mathrm{R}G(Y) \right), \]
	where $X\in\Obj\left(\cate{D}(\mathcal{A})\right)$ and
	$Y\in\Obj\left(\cate{D}(\mathcal{A}')\right)$.
\end{theorem}
% segment-id: o014.li2.chapter4.unbounded-derived-functors.pf008
\begin{proof}
	The argument is the same as for
	Theorem~\ref{prop:RHom-adjoint-bounded}. The only difference is that
	injective (respectively, projective) resolutions are replaced by
	K-injective (respectively, K-projective) resolutions, and
	Theorem~\ref{prop:K-resolution-homotopy} replaces
	Theorem~\ref{prop:resolution-homotopy} in the discussion of uniqueness of
	resolutions.
\end{proof}

% segment-id: o014.li2.chapter4.unbounded-derived-functors.r001
\begin{remark}
	Applying $\Hm^0$ to both sides of Theorem~\ref{prop:RHom-adjoint} gives the
	adjunction
% segment-id: o014.li2.chapter4.unbounded-derived-functors.q009
	\[ \Hom_{\cate{D}(\mathcal{A'})}\left(\mathrm{L}F(X), Y\right) \simeq \Hom_{\cate{D}(\mathcal{A})}\left(X, \mathrm{R}G(Y) \right). \]

	Since derived functors obtained through resolutions are always absolute Kan
	extensions in the sense of Definition~\ref{def:absolute-Kan}
	(Proposition~\ref{prop:localization-functor-abs}),
	Theorem~\ref{prop:Quillen-adjunction-gen} also says that
	$(\mathrm{L}F,\mathrm{R}G)$ is an adjoint pair. The two approaches give the
	same adjunction isomorphism. Why? It suffices to consider the case in which
	$X$ is K-projective and $Y$ is K-injective. In this case, the adjunction
	isomorphism of Theorem~\ref{prop:RHom-adjoint} is realized concretely by the
	unit $\eta$ and counit $\varepsilon$ at the level of $\cate{K}(\cdot)$ as
% segment-id: o014.li2.chapter4.unbounded-derived-functors.q010
	\begin{align*}
		\Hom_{\cate{K}(\mathcal{A}')}\left( \cate{K}F(X), Y\right) \simeq \Hom_{\cate{K}(\mathcal{A})}\left( X, \cate{K}G(Y) \right).
	\end{align*}
	On the other hand, Theorem~\ref{prop:Quillen-adjunction-gen} determines
	$\underline{\eta}$ and $\underline{\varepsilon}$ at the level of derived
	categories. The characterization given there suffices to show that they are
	compatible with $\eta$ and $\varepsilon$. The details are left to the
	interested reader.
\end{remark}

% segment-id: o014.li2.chapter4.unbounded-derived-functors.p011
Finally, another set of sufficient conditions for the existence of unbounded
derived functors is based on finite dimension
(Definition--Proposition~\ref{def:dim-functor}). Although the proof technique
does not go beyond the scope of this book, we state the result without proof
to avoid a digression; see {\cite[Tags 07K6, 07K7]{stacks}}. The starting
point is the construction of bounded derived functors in
Theorem~\ref{prop:derived-via-resolution}.

% segment-id: o014.li2.chapter4.unbounded-derived-functors.pr003
\begin{proposition}\label{prop:unbounded-derived-fd}
% segment-id: o014.li2.chapter4.unbounded-derived-functors.x004
	\index{dimension}
	Let $F:\mathcal{A}\to\mathcal{A}'$ be a left-exact (respectively,
	right-exact) functor between abelian categories. Assume that
% segment-id: o014.li2.chapter4.unbounded-derived-functors.l004
	\begin{compactitem}
% segment-id: o014.li2.chapter4.unbounded-derived-functors.i009
		\item relative to $F$, there is a full additive subcategory
		$\mathcal{A}^\flat$ of $\mathcal{A}$ of type I (respectively, type P),
		as in Definition--Proposition~\ref{prop:F-injective-criterion};
% segment-id: o014.li2.chapter4.unbounded-derived-functors.i010
		\item the corresponding bounded derived functor
		${}^+\mathrm{R}F$ (respectively, ${}^-\mathrm{L}F$) has finite
		dimension.
	\end{compactitem}
	Then $\mathrm{R}F$ (respectively, $\cate{L}F$):
	$\cate{D}(\mathcal{A})\to\cate{D}(\mathcal{A}')$ exists. Moreover,
	$\cate{K}(\mathcal{A}^\flat)$ is an $F$-injective (respectively,
	$F$-projective) subcategory of $\cate{K}(\mathcal{A})$.
\end{proposition}

% segment-id: o014.li2.chapter4.unbounded-derived-functors.ex001
\begin{example}\label{eg:Rlim-unbounded}
	Suppose that $\mathcal{A}$ has exact countable products and enough
	injective objects. Take the left-exact functor $F$ to be
	$\varprojlim:\cate{InvSys}(\mathcal{A})\to\mathcal{A}$.
	Theorem~\ref{prop:Rlim-as-holim} (iii) shows that the bounded version of
	$\mathrm{R}\lim$ has dimension at most $1$, while the subcategory
	$\mathcal{R}\subset\cate{InvSys}(\mathcal{A})$ is of type I relative to
	$\varprojlim$. The preceding result therefore gives the unbounded derived functor
% segment-id: o014.li2.chapter4.unbounded-derived-functors.q011
	\[ \mathrm{R}\lim: \cate{D}(\cate{InvSys}(\mathcal{A})) \to \cate{D}(\mathcal{A}). \]

	As in Theorem~\ref{prop:Rlim-as-holim} (ii), for every object $X$ of
	$\cate{C}(\cate{InvSys}(\mathcal{A}))$ there is a distinguished triangle
	in $\cate{D}(\mathcal{A})$
% segment-id: o014.li2.chapter4.unbounded-derived-functors.q012
	\[ \mathrm{R}\lim(X) \to \prod_k X_k \xrightarrow{\Delta_X} \prod_k X_k \xrightarrow{+1} . \]
	In particular, $\mathrm{R}\lim(X)\simeq\holim(X)$. By
	Proposition~\ref{prop:unbounded-derived-fd}, the original proof applies with
	no other changes: simply replace every $\cate{C}^+(\cdots)$ by
	$\cate{C}(\cdots)$.
\end{example}
